% INFO6007 — Project Management in IT — Exam Cheatsheet (S2 2026)
% One A4 sheet, double-sided. Run-in paragraphs at 5.7pt. Page 1: W1–W6 + question bank Q1–Q24. Page 2: W7–W12 + question bank + practice exam full answers.
% Compile: pdflatex INFO6007_cheatsheet.tex  (must stay at 2 pages)
\documentclass[a4paper,landscape]{article}
\usepackage{xcolor}
\usepackage[top=5mm,bottom=5mm,left=6mm,right=6mm]{geometry}
\usepackage{iftex}
\ifPDFTeX
  \usepackage[utf8]{inputenc}
  \usepackage[T1]{fontenc}
\fi
\usepackage{lmodern}
\usepackage{amsmath}
\usepackage{amssymb}
\usepackage{microtype}
\usepackage{multicol}
\usepackage{enumitem}
\setlist[itemize]{leftmargin=1.6mm,topsep=0pt,itemsep=0pt,parsep=0pt,label=\textbullet}
\setlength{\columnsep}{5pt}\setlength{\smallskipamount}{0.6pt}
\setlength{\parindent}{0pt}
\binoppenalty=50
\relpenalty=50
\emergencystretch=1.5em
\renewcommand{\familydefault}{\sfdefault}
\pagestyle{empty}

% ---- colour commands (same macros as the COMP5270 sheet) ----
\newcommand{\hlhighlight}[1]{\begingroup\setlength{\fboxsep}{0pt}\colorbox{yellow}{\strut\textbf{#1}}\endgroup}
\newcommand{\sectitle}[1]{\textbf{#1}}
\newcommand{\concept}[1]{\textbf{#1}}
\newcommand{\lbl}[1]{\textbf{#1}}
\definecolor{keyhlcolor}{HTML}{FFEB7A}
\newcommand{\kh}[1]{{\setlength{\fboxsep}{0.5pt}\colorbox{keyhlcolor}{#1}}}
\newcommand{\ra}{$\rightarrow$}

\newcommand{\PONE}{\fontsize{5.7pt}{6.05pt}\selectfont}
\newcommand{\PTWO}{\fontsize{5.7pt}{6.05pt}\selectfont}

\begin{document}
\PONE
\sloppy

%=====================================================================
%  PAGE 1 — WEEKS 1–6
%=====================================================================
\begin{multicols}{3}
\PONE

\par\smallskip\noindent\hlhighlight{W1 · Projects, success, triple constraint}\ \concept{Project} (PMBOK): \kh{temporary effort} to create a \kh{unique product, service or result}; clear start/end, specific objectives, constrained by scope/time/cost/resources. \concept{Operations} = ongoing, repetitive (bug fixes, user queries, BAU support) \ra\ ITSM. Weekly new features after go-live with no end = scope creep / should be a new project or product backlog with own funding. \textbullet\ \concept{PM} = applying knowledge, skills, tools \& techniques to project activities to meet requirements. Elements: scope, time, cost, quality, risk, communication, resources. \textbullet\ \concept{Why IT projects are hard:} evolving requirements, intangible outcomes, hidden integration complexity, decisions before full info, rapid tech change. \lbl{Failure patterns:} volatile reqs, underestimated integration, optimism bias, business--IT misalignment, weak governance/decision rights, low stakeholder involvement. \textbullet\ \concept{Triple constraint:} scope $\uparrow$ \ra\ time and/or cost $\uparrow$; time $\downarrow$ \ra\ cost $\uparrow$ or scope $\downarrow$; budget $\downarrow$ \ra\ scope $\downarrow$ or time $\uparrow$ or quality/risk hit. Quality sits in the middle: \kh{never let quality be the hidden constraint} that absorbs a change. Options: deliver smaller first release, stage new features, or formally revise baseline. \textbullet\ \concept{Success} $\ne$ only on time/on budget: + quality, stakeholder satisfaction, business value / benefits realisation, user adoption, strategic alignment, team performance, sustainability. Misconceptions: ``on time \& budget = success''; ``all success is immediate''. Set \lbl{measurable success criteria at initiation} (e.g.\ adoption $\ge$70\% in 3 months, processing time $-$50\%), agreed by sponsor + key users, recorded in charter; benefits owner measures after closure. \textbullet\ \concept{Project / Program / Portfolio:} project = tactical, one output (PM); \lbl{program} = related projects coordinated for benefits not available separately (program mgr; interdependencies); \lbl{portfolio} = projects + programs selected for strategic goals (portfolio mgr / PMO; investment \& prioritisation). \textbullet\ \concept{Sunk cost:} money already spent is irrelevant to go/no-go; decide on \kh{future costs vs future benefits}. Terminated project still needs closure: accept/handover completed work, close contracts, release resources, archive, \lbl{lessons learned}.\par

\par\smallskip\noindent\hlhighlight{W2 · Lifecycle \& methodologies}\ \concept{Lifecycle:} Initiation (charter, approved) \ra\ Planning (PM plan) \ra\ Execution (deliverables) \ra\ Monitoring \& Controlling (status reports, change control; runs alongside execution) \ra\ Closure (final report, lessons learned, sign-off). Skipping initiation (building prototypes before charter) \ra\ no authorised budget, unknown stakeholders, no agreed scope. \textbullet\ \concept{Waterfall:} linear; Reqs \ra\ Design \ra\ Code \ra\ Test \ra\ Deploy \ra\ Maintain. + simple, clear milestones \& documentation, fixed scope, compliance/audit, outsourced handoffs, fixed-price contracts. $-$ inflexible, late testing finds problems late, costly change. \lbl{Planning} upfront \& detailed; \lbl{stakeholders} at start and end; \lbl{change} via formal CR, expensive. \textbullet\ \concept{Agile} = mindset; iterative, incremental, working software each sprint; customer collaboration \& responding to change. \lbl{Planning} rolling per sprint; \lbl{stakeholders} continuous (reviews/demos); \lbl{change} welcomed via backlog reprioritisation. Risk: scope drift, weak documentation, stakeholder fatigue. \textbullet\ \concept{Scrum:} sprints 1--4 wk; \lbl{roles} Product Owner (backlog, priority), Scrum Master (coach, removes blockers), Dev team (cross-functional); \lbl{artefacts} product backlog, sprint backlog, increment; \lbl{events} sprint planning, daily stand-up (15 min), review (demo), retrospective. \textbullet\ \concept{Kanban:} continuous flow, board To Do/Doing/Done, \kh{WIP limits}, cycle time, pull when capacity; no fixed roles; change anytime \ra\ support/ops ticket streams. \concept{DevOps:} Dev+Ops; CI/CD, IaC, automated testing \& monitoring, faster reliable releases (+DevSecOps, MLOps, FinOps...). \textbullet\ \concept{Hybrid (water-scrum-fall):} Waterfall for compliance, requirements sign-off, procurement milestones; Agile for build \& feedback. Fits regulated health/gov apps. Program with Scrum team + Kanban ops + Waterfall hardware rollout = hybrid; coordinate via shared roadmap/milestones, integration manager, program-level board \& sync meetings. \textbullet\ \lbl{Pick:} fixed legislation, fixed-price, audit \ra\ Waterfall; uncertain features, user feedback, start-up \ra\ Agile/Scrum; staged release + later features (e.g.\ uni app) \ra\ Agile (core first).\par

\par\smallskip\noindent\hlhighlight{W2 · Requirements, initiation, charter, stakeholders}\ \concept{Why requirements matter:} product meets needs, reduces scope creep \& rework (defects from undocumented expectations), basis for design/test/estimates. \concept{Techniques:} interviews (depth, 1:1), surveys (many users fast), workshops/JAD (reconcile conflicts), observation (real workflow, busy staff), document analysis (existing manuals), prototyping (refine vague needs), focus groups. Busy staff + 5 clinics + outdated manual \ra\ observation + short interviews + one workshop + prototype tested with elderly/limited-English users. \textbullet\ \concept{Functional} = what system does (view usage, generate report). \concept{Non-functional} = how well: performance, availability, security, usability, scalability, accessibility. \kh{Good req: specific, measurable, testable, unambiguous} --- ``fast and user friendly'' is bad \ra\ ``dashboard loads $\le$3 s with 25{,}000 concurrent users''; ``99.95\% monthly availability''; ``new user books in $\le$3 clicks, SUS $\ge$75''. \lbl{QB Q11:} view usage (a) \& admin report (d) = functional; 3 s @25k users (b) = performance/scalability; encryption (c) = security; 99.95\% (e) = availability/reliability. Verify (b) by load/stress test at 25k simulated users (evidence: p95 response report); (e) by uptime monitoring over a month (evidence: SLA availability report). Costliest usually (e) or (b) at peak weather events \ra\ redundancy/auto-scaling; balance by phasing, cloud elasticity, negotiating target (99.9\%), MoSCoW. \textbullet\ \concept{Conflicts} (finance 5-step vs ops 1-step): workshop with both, link to business rules/risk, sponsor decides, \lbl{MoSCoW} prioritisation, configurable rule (e.g.\ 1 step $<$\$X, multi-step above). \textbullet\ \concept{Traceability (RTM):} links each requirement \ra\ source \ra\ design \ra\ test case \ra\ acceptance; spots gold-plating \& missing reqs, supports scope control, impact analysis \& test coverage. \lbl{Baseline} = signed-off requirements; validate with walkthroughs/prototypes/UAT criteria. \textbullet\ \concept{Initiation activities:} business need, objectives, high-level scope, stakeholders, appoint PM, \lbl{project charter}, high-level budget/timeline, feasibility \& initial risks, strategic alignment. Tools: expert judgement, business case, SWOT, workshops. \textbullet\ \concept{Project charter} = formally authorises project, gives PM authority to use resources. Contains: purpose/business case, objectives \& success criteria, high-level scope/deliverables, milestones, budget, key stakeholders, sponsor, risks/assumptions/constraints, approval. \textbullet\ \concept{Stakeholder} = anyone who affects or is affected. Process: identify \ra\ analyse (power/interest/influence) \ra\ prioritise \ra\ engage \ra\ monitor. \kh{Power/Interest grid:} high P + high I = \lbl{manage closely}; high P + low I = \lbl{keep satisfied}; low P + high I = \lbl{keep informed}; low P + low I = \lbl{monitor}.\par

\par\smallskip\noindent\hlhighlight{W3 · Scope management (6 processes)}\ \concept{Plan scope mgmt} \ra\ \concept{Collect requirements} (reqs doc + RTM) \ra\ \concept{Define scope} (scope statement) \ra\ \concept{Create WBS} \ra\ \concept{Validate scope} \ra\ \concept{Control scope}. Scope mgmt = all the work and \emph{only} the work required. \textbullet\ \concept{Scope statement:} product description, deliverables, \lbl{inclusions/exclusions}, acceptance criteria, assumptions, constraints. A one-paragraph statement with no acceptance criteria \ra\ ``doesn't do what we expected'' dispute. \textbullet\ \concept{WBS:} \kh{deliverable-oriented hierarchical decomposition} of the total scope; levels \ra\ \lbl{work packages} (lowest; one owner, estimate, testable output). Backbone for activity list, schedule, cost estimates, resources, risk ID, responsibility. Rules: each unit of work in one place only (\emph{no double-counted data migration}); parent = sum of children (100\% rule); one accountable owner per item; team builds it; every item in \lbl{WBS dictionary} (description, owner, deliverable, acceptance, estimate). Approaches: guidelines, analogy, top-down, bottom-up, mind-map. \lbl{Phase-based} WBS (Reqs/Design/Build/Test) follows lifecycle; \lbl{deliverable-based} (Portal/Integration/Training...) makes ownership \& completeness clearer. \concept{Scope baseline} = scope statement + WBS + WBS dictionary \ra\ protects against ``not our work'' disputes. \textbullet\ \concept{Validate scope} = \kh{formal acceptance of deliverables by customer/sponsor} (inspection, sign-off; outputs accepted deliverables, CRs). \concept{Control scope} = monitor scope vs baseline, manage changes via variance analysis \& \lbl{integrated change control}. (Validate = acceptance/\emph{right thing}; Control = managing changes; QC = correctness.) \textbullet\ \concept{Scope creep} = \kh{uncontrolled growth of scope} without adjusting time, cost, resources (e.g.\ devs add features w/o impact assessment or approval). Consequences: delay, budget overrun, rushed testing \ra\ lower quality, burnout, stakeholder conflict. Causes: vague scope, no change control, weak stakeholder engagement, gold-plating. Controls: \lbl{formal change control} (log CR \ra\ impact on scope/time/cost/quality/risk \ra\ CCB/sponsor approve/defer/reject \ra\ update baseline/WBS/RTM \ra\ communicate), clear scope statement, RTM, sign-offs per stage, regular scope reviews, defer to later release/backlog.\par

\par\smallskip\noindent\hlhighlight{W3--4 · Schedule management}\ \concept{6 processes:} plan schedule mgmt \ra\ define activities (from work packages) \ra\ sequence activities \ra\ estimate durations \ra\ develop schedule (CPM, baseline) \ra\ control schedule (SPI/SV, compression). \textbullet\ \concept{Dependencies:} \lbl{mandatory} (hard logic, inherent: build before test; licence before deploy), \lbl{discretionary} (soft logic, preference/best practice: design after procurement --- can be changed to fast-track, raising risk), \lbl{external} (outside team control: vendor, gov agency, licence), \lbl{internal} (within team). A dep can be mandatory+external. \textbullet\ \concept{Relationships (PDM):} FS (most common), SS, FF, SF (rare). \concept{Lead} = successor starts before predecessor finishes (speeds up; e.g.\ testers write scripts while C in progress = SS or FS with lead); \concept{lag} = wait after predecessor (slows). Gantt: overlapping bars/bar shifted left = lead. \textbullet\ \concept{Network diagrams:} AON (nodes = activities, most used), AOA (arrows = activities, may need dummies), PDM. \concept{Gantt} = bars vs time, best for communicating timeline \& progress; \lbl{network diagram} shows logic/dependencies \& critical path. \textbullet\ \concept{CPM:} \kh{critical path = longest path} = shortest project duration; \kh{zero float}; delay on it delays project. Forward pass ES/EF (EF=ES+d, take max of preds), backward pass LS/LF (take min of succs); \kh{Total float = LS $-$ ES = LF $-$ EF}. Delay $\le$ float \ra\ no end-date change; delay $>$ float \ra\ project slips by (delay $-$ float) and path may become critical. \lbl{QB Q16:} A3, B5(A), C4(A), D6(B), E2(C), F3(D,E): paths A-B-D-F = 17, A-C-E-F = 12 \ra\ \kh{CP A-B-D-F, 17 d}; float C = E = 5; C late 3 d \ra\ no change. \textbullet\ \concept{Estimating:} \lbl{analogous} (top-down, historical, fast, less accurate, early phase), \lbl{parametric} (unit rate $\times$ units), \lbl{bottom-up} (work packages rolled up; most accurate, slow), \lbl{three-point/PERT}, expert judgement; rolling-wave for later phases. \textbullet\ \concept{PERT:} \kh{$t_E=(O+4M+P)/6$}, $\sigma=(P-O)/6$. Practice: O=4, M=7, P=16 \ra\ $(4+28+16)/6=$ \kh{8 days}, $\sigma=2$. Useful: weights uncertainty, exposes risk range. Limit: still an estimate --- skewed toward pessimistic (16), only $\sim$50\% chance of $\le$8 d; commit to $t_E+\sigma$ to $+2\sigma$ (10--12 d) or add buffer; depends on quality of O/M/P guesses. \textbullet\ \concept{Compression:} \kh{crashing} = add resources to \lbl{critical} activities (cost $\uparrow$; pick cheapest cost/day; crashing non-critical activity with float is waste). \kh{Fast tracking} = run sequential activities in parallel (risk/rework $\uparrow$, cost ~same). Scope reduction (MVP) or accept late date are \emph{not} compression. \lbl{Brooks' law:} adding contractors late \ra\ ramp-up (1 wk learning) + communication overhead ($n(n-1)/2$ channels) may slow further.\par

\par\smallskip\noindent\hlhighlight{W4 · Cost management \& EVM}\ \concept{4 processes:} plan cost mgmt \ra\ estimate costs \ra\ determine budget \ra\ control costs. Cost types: tangible/intangible, \lbl{direct} (devs, licences) / \lbl{indirect} (overhead, office), \lbl{sunk} (ignore). \textbullet\ \concept{Reserves:} \lbl{contingency} = known-unknowns (identified risks), \emph{inside} cost baseline, PM uses; \lbl{management} = unknown-unknowns, \emph{outside} baseline, needs sponsor approval. \kh{Cost baseline = estimates + contingency}; \kh{Budget = baseline + management reserve}. Baseline = approved time-phased budget (= cumulative PV, S-curve). \textbullet\ \concept{Estimating pitfalls:} underestimating integration/testing/maintenance, hidden costs (training, licences, cloud consumption, security, vendor support), optimism, currency. Benchmarks (\$/user migrated, sprint cost). \textbullet\ \concept{EVM:} PV = budgeted cost of work \emph{planned}; EV = budgeted cost of work \emph{performed} (\%complete $\times$ BAC); AC = actual cost. \textbullet\ \kh{CV = EV $-$ AC} (+ under budget); \kh{SV = EV $-$ PV} (+ ahead); \kh{CPI = EV/AC}; \kh{SPI = EV/PV} ($<1$ bad); \kh{EAC = BAC/CPI} (trend continues); EAC = AC + (BAC $-$ EV) (variance atypical); EAC = AC + (BAC$-$EV)/(CPI$\times$SPI); ETC = EAC $-$ AC; \kh{VAC = BAC $-$ EAC}; TCPI = (BAC$-$EV)/(BAC$-$AC). \textbullet\ \lbl{Practice Q4a:} PV 300k, EV 240k, AC 320k \ra\ \kh{SPI 0.80} (80\% of planned progress, behind), \kh{CPI 0.75} (\$0.75 value per \$1, over budget); SV $-$60k, CV $-$80k. \lbl{QB Q22:} BAC 500k, PV 270k, EV 200k, AC 230k \ra\ SV $-$70k, CV $-$30k, SPI 0.74, CPI 0.87, EAC 575k, VAC $-$75k. \lbl{QB Q23:} A (SPI 1.10, CPI 0.80) ahead but overspending --- likely crashing/overtime; B (SPI 0.80, CPI 1.10) behind but cheap --- under-resourced; concern depends on whether B's slip is on critical path / deadline fixed vs budget fixed. \textbullet\ \concept{Corrective actions:} root cause (integration, change requests, resource churn, vendor delay, cloud sprawl), scope triage/defer low-value, reallocate senior devs to critical path, enforce change control, renegotiate vendor, reforecast; compression only if schedule matters more than cost (crashing worsens CPI). \textbullet\ \concept{Agile cost:} fund per sprint/release (time \& cost fixed, \lbl{scope flexible}), burn-up/down, velocity; EVM hard (PV unstable, story points $\ne$ \$) \ra\ EBV, map points to cost via burn rate, re-forecast EAC each iteration; still keep contingency.\par

\par\smallskip\noindent\hlhighlight{W5 · Quality management}\ \concept{Quality} = \kh{degree to which inherent characteristics} fulfil requirements: conformance to requirements + fitness for use + customer satisfaction. Dimensions: functionality, reliability, performance, usability, maintainability, security. Principles: customer focus, \lbl{prevention over inspection}, continuous improvement, fact-based decisions. \textbullet\ \concept{3 processes:} \lbl{Plan quality} (standards, metrics, roles, QA/QC activities \ra\ QMP) \ra\ \lbl{Manage quality = QA} (\emph{process}, preventive: process audits, peer/code reviews, design walkthroughs, retrospectives, coding standards, static analysis, PDCA/DMAIC) \ra\ \lbl{Control quality = QC} (\emph{product}, reactive: testing, inspection, defect logging, compare to metrics, accept/conditional/reject). \kh{QA = right process; QC = right output.} \lbl{Verification} = built it right (spec); \lbl{validation} = built the right thing (user needs). \textbullet\ \concept{Testing:} unit, integration, system, \lbl{UAT}, performance/load (peak enrolment case: passed all planned tests but no load test with realistic concurrency \ra\ planning failure), security, regression; alpha (controlled env) / beta (real users). \textbullet\ \concept{Tools:} cost-benefit, \lbl{cause-and-effect/Ishikawa} (root causes: people, process, tools, environment, data), \lbl{control chart} (process stable? UCL/LCL, trends), checksheet (tally defects), flowchart (process gaps), scatter (correlation), histogram (frequency), \lbl{Pareto} (\kh{80/20}: fix vital few causes first). \textbullet\ \concept{Standards:} ISO 9001 (QMS), ISO/IEC 27001 (infosec), CMMI 1 Initial \ra\ 2 Managed \ra\ 3 Defined \ra\ 4 Quant.\ managed \ra\ 5 Optimising; ITIL; Lean (cut waste), Six Sigma (cut variation, DMAIC). \textbullet\ \concept{Cost of quality:} \kh{CoC = prevention + appraisal} (training, reviews, automated tests, audits); \kh{CoNC = internal failure + external failure} (rework, bug fixes before release; downtime, complaints, penalties after). e.g.\ skipping \$2k prevention \ra\ risk \$20k external failure \ra\ invest in prevention. More testers at end = QC/appraisal; standards + reviews + automated checks early = QA/prevention (cheaper). \textbullet\ \concept{Deming:} build quality into process, cease dependence on inspection, drive out fear, break barriers, training, constancy of purpose, everyone's job. \concept{Juran:} top-management commitment; 10 steps (awareness, goals, organise, resources, train, act, report, recognise, communicate, maintain momentum). \textbullet\ \lbl{QB Q6:} 1750 scans, 105 fail \ra\ pass $=1645/1750=$ \kh{94\% $<$ 97\%} \ra\ reject batch. QA: agreed scanning standards + supplier process audit/calibration; QC: sample inspection of each batch vs resolution spec.\par

\par\smallskip\noindent\hlhighlight{W6 · Resource management, people \& teams}\ \concept{6 processes:} plan resources \ra\ estimate activity resources \ra\ acquire (internal, external hire/contractor, procurement, multi-source) \ra\ \lbl{develop team} \ra\ \lbl{manage team} \ra\ control resources. Human + physical (servers, cloud, licences, devices). \textbullet\ \concept{RACI:} \lbl{R}esponsible (does work), \lbl{A}ccountable (\kh{exactly one}, approves), \lbl{C}onsulted (two-way, before), \lbl{I}nformed (one-way, after). e.g.\ ``approve design change'': A = PM/design authority; R = lead architect; C = senior engineers, contractors' lead, security; I = team, sponsor \ra\ removes who-approves confusion. \concept{RBS} = resources by category (Human/Equipment/Software/Materials \ra\ roles \ra\ named items). \textbullet\ \concept{Resource histogram} = allocation per period, shows over-allocation. \concept{Loading} = allocated/available hours (Dev1 50/40 = 125\% over; Dev2 30/40 under). \kh{Levelling} = resolve over-allocation by delaying/splitting/reassigning tasks --- \lbl{may extend finish date / change critical path}. \kh{Smoothing} = adjust only within float, end date fixed. Lead dev 160\% on critical path + idle testers: reassign suitable tasks (testers do test prep, code review), add contractor, overtime short-term, split tasks; risks = skill gaps, burnout, ramp-up. \textbullet\ \concept{Leadership styles:} autocratic/directive (fast, crises, \emph{new inexperienced staff need clear direction}), democratic/participative (engagement, slower), laissez-faire (\emph{experienced experts want autonomy}), transformational, transactional, servant, \lbl{situational/adaptive} --- \kh{match style to team maturity}: directive for novices + coaching, delegating for experts (one-size directive = demotivates seniors, stifles creativity). \textbullet\ \concept{Motivation:} intrinsic vs extrinsic; \lbl{Maslow} (physiological \ra\ safety \ra\ belonging \ra\ esteem \ra\ self-actualisation); \lbl{Herzberg} hygiene (salary, policy, conditions, security --- absence dissatisfies) vs motivators (achievement, recognition, responsibility, growth --- create motivation); Covey 7 habits (win-win, seek first to understand, synergise). \textbullet\ \concept{Tuckman:} \kh{Forming \ra\ Storming \ra\ Norming} \kh{\ra\ Performing \ra\ Adjourning}. \lbl{New members join a performing team \ra\ team regresses to forming/storming} (disputes over working practices, slower delivery) --- sponsor's ``experienced team needs no support'' is wrong. Intervention: onboarding + team charter/working agreements workshop, pair new devs with mentors, clarify RACI, retrospective; PM as facilitator/coach. \textbullet\ \concept{Conflict modes:} competing (win-lose; urgent, critical), \lbl{collaborating} (win-win; complex, important, time available --- best for technical-direction clash), compromising (partial; time pressure), avoiding (trivial/emotions high), accommodating (preserve relationship). Lencioni dysfunctions: absence of trust, fear of conflict, lack of commitment, avoidance of accountability, inattention to results. \textbullet\ \concept{Agile resources:} capacity planning per sprint, cross-functional self-managing teams, watch WIP/burnout, rebalance each sprint.\par

\par\smallskip\noindent\hlhighlight{Question bank · W1--W6 answer skeletons}\ \lbl{Q1 telehealth ``never finished'':} weekly new features = new scope with no end \ra\ should be a new project/release or product backlog with own funding; answering queries + minor bug fixes = \kh{operations/BAU} \ra\ hand to service desk/ITSM. Creep: features keep flowing in without CRs, no closure criteria \ra\ no point where sponsor signs off, team never released, finance cannot measure ROI because baseline \& end date blurred. Fix at initiation: measurable success criteria (uptake, wait time, cost per consult), defined deliverables + acceptance, closure trigger, transition plan to operations, benefits owner reviews 6--12 months later. \textbullet\ \lbl{Q2 fixed budget + legal deadline, +30\% effort:} scope $\uparrow$ 30\% with time \& cost fixed \ra\ only scope or quality can move; quality must not be the hidden victim (regulatory tool). Options: (1) \kh{defer analytics to phase 2} after deadline (+ meets law, $-$ client waits); (2) extra budget/contractors for parallel build (+ all scope, $-$ cost, Brooks ramp-up, integration risk); (3) descope low-value must-nots elsewhere to fit a minimal analytics slice (MoSCoW). Rec: \lbl{scope is the flexible constraint} (legal date immovable, budget fixed) \ra\ CR to CCB, defer module, record decision \& update baseline. \textbullet\ \lbl{Q3 A vs B:} A = triple-constraint success but 20\% adoption \ra\ no benefits, value failure; B = failed time/cost (+20\%, 1 month late) but halved processing time, sponsor satisfied \ra\ business success. Success criteria set at initiation, agreed by \kh{sponsor + key users/benefits owner}, recorded in charter/business case, mix of delivery + adoption + benefit KPIs with targets \& measurement dates. \textbullet\ \lbl{Q4 prototypes before charter:} skipped \kh{initiation} \ra\ missing charter, business case, stakeholder register, high-level scope/budget, PM authority. Consequences: unauthorised spend (sponsor questions budget), scope based on guesses \ra\ rework, unidentified risks, excluded stakeholders resist. Recover: pause build, draft charter + business case, stakeholder analysis \& engage the two missed stakeholders, treat prototype as requirements input, get sponsor sign-off, then plan. \textbullet\ \lbl{Q5 phase gate, +25\% cost \& supplier gone:} work already spent = \kh{sunk cost}, irrelevant; decide on remaining cost (now +25\%) vs remaining benefit + new supplier risk; gate exists exactly to kill/re-scope. If terminated: formal closure --- accept \& archive completed designs, close/settle contracts, release staff, final report, lessons learned (estimating, supplier due diligence) \ra\ reused by PMO so next estimate/supplier choice is better. \textbullet\ \lbl{Q7:} licensing system \ra\ \kh{Waterfall} (fixed legislation scope, fixed-price contract, audit docs); start-up app \ra\ \kh{Agile/Scrum} (features from feedback, fast releases). Change: formal CR vs backlog reprioritisation; docs: comprehensive vs just enough; customer: start/end sign-off vs every sprint review. Risk W: late discovery of usability issues \ra\ prototype + staged UAT; risk A: scope drift/runway burn \ra\ product goal, MVP, sprint budget caps. \textbullet\ \lbl{Q8 program:} dev = Scrum (2-wk sprints, roles, timebox), ops = Kanban (continuous flow, WIP limits, no sprints), hardware = Waterfall (contract milestones) \ra\ \kh{hybrid} = different methods under one program. Challenges: misaligned cadences/dependencies (sprint needs hardware) \ra\ integrated roadmap + milestone sync; inconsistent reporting \ra\ program board, common KPIs, scrum-of-scrums/integration manager. \textbullet\ \lbl{Q9/Q10 requirements:} root causes = no elicitation with real users, undocumented expectations, no sign-off/baseline, no RTM, no validation \ra\ 30\% defects. Validate via walkthroughs, prototypes, UAT acceptance criteria per req; baseline = signed-off set under change control. RTM links req \ra\ design \ra\ test \ra\ acceptance: every req tested, no orphan features, impact analysis for CRs. \textbullet\ \lbl{Q12 ``secure and easy to use'':} rewrite \ra\ FR: role-based access (only treating clinicians view records), audit log of every access; NFR: MFA for all logins, AES-256 at rest/TLS 1.2+, record opens $\le$2 s for 95\% of requests, new clinician completes task in $\le$3 clicks/SUS $\ge$75. Each req gets acceptance criterion + test (pen test, load test, usability test). Vague NFRs \ra\ built to wrong target, late defects, expensive rework, audit failure, low adoption. \textbullet\ \lbl{Q13 portal rejected:} \kh{Validate Scope} = formal acceptance of deliverables by client; \kh{Control Scope} = managing changes to baseline. Scope statement needed: deliverables, inclusions/exclusions, acceptance criteria, assumptions, constraints. Resolve: compare delivered vs documented scope, joint review, agree acceptance criteria now, fix true defects, raise CRs for new expectations; future: interim reviews/demos + sign-off per deliverable. \textbullet\ \lbl{Q14 ERP flat list:} WBS L1 ERP system \ra\ L2 Finance module, HR module, \kh{Data migration (one owner)}, Integration, Testing (unit/system/UAT/perf as separate work packages, not one 120-day item), Training (materials, sessions, \kh{owner named}), Deployment, PM \ra\ L3 work packages. WBS dictionary defines each package (scope, owner, estimate) \ra\ no double counting; RAM/RACI gives exactly one A per package \ra\ training owned. \textbullet\ \lbl{Q15:} WBS + dictionary + scope baseline = documented, agreed 100\% of scope \ra\ site documentation must appear as a package with owner \& contract reference; WBS \ra\ activities (decompose packages) \ra\ sequence/durations \ra\ schedule; package estimates roll up to cost baseline. \textbullet\ \lbl{Q17:} licence before deployment = \kh{mandatory + external} (no flexibility; track agency, buffer); design after procurement = \kh{discretionary + internal} \ra\ remove to run in parallel = fast tracking \ra\ shorter critical path but rework risk if procurement changes specs. Network diagram = logic, dependencies, critical path, float (planning/analysis); Gantt = timeline, durations, progress (communication). \textbullet\ \lbl{Q18 3 wks behind:} contractors on X: 1-wk ramp-up, more channels $n(n-1)/2$, mentoring slows experts \ra\ \kh{Brooks's law}, may not save time. Crash X (critical, 3 d for \$6k = \$2k/day); \kh{do not crash Y} (8 d float, no effect on end date, \$1k wasted). Fast-track non-dependent work (e.g.\ test prep, docs in parallel) for remaining gap, watch rework risk. \textbullet\ \lbl{Q19 go-live 4 wks earlier:} build \& test in parallel = \kh{fast tracking} (compression; more risk/rework, no extra cost); MVP = \kh{scope reduction} (not compression; less value at launch); accept late launch = no compression (misses marketing event). Compare cost, risk, quality, stakeholders \ra\ rec MVP for event + rest in next release (or MVP + limited fast-track); communicate via sponsor briefing + updated plan, CR approved, rebaseline scope/schedule. \textbullet\ \lbl{Q20/Q21:} see W6 levelling vs smoothing (smoothing cannot fix 160\% on critical path \ra\ levelling or reassignment); remote team: RACI per decision + team charter (hours overlap, channels), collaborating for engineer clash (technical spike/prototype + data), compromise if deadline, PM escalates to design authority if stuck. \textbullet\ \lbl{Q24 defects climbing:} more testers at end = QC/appraisal (finds defects late, costly internal failure); standards + reviews + automated checks = QA/prevention. CoQ: \kh{prevention $<$ appraisal $<$ internal $<$ external failure} cost \ra\ invest early. Tools: \lbl{Pareto} on defect categories (vital few modules), \lbl{fishbone} for root causes, \lbl{control chart} of defect rate per build to see if process stable after fix.\par

\end{multicols}

\newpage
%=====================================================================
%  PAGE 2 — WEEKS 7–12 + EXAM TOOLKIT
%=====================================================================
\begin{multicols}{3}
\PTWO

\par\smallskip\noindent\hlhighlight{W7 · Risk management}\ \concept{Risk} = \kh{uncertain event/condition} that, if it occurs, has a \kh{+ or $-$ effect} on objectives. Threats vs opportunities. Categories: market, financial, technology, people, structure/process. \lbl{Risk vs issue:} risk = may happen (future); issue = has happened. \textbullet\ \concept{Risk appetite} = uncertainty an org is willing to accept (startup high vs government/museum/bank low); culture risk-averse / neutral / seeking; \lbl{utility}: averse prefers certain lower payoff, neutral uses EMV only, seeking gambles. Appetite must match culture. \textbullet\ \concept{6 processes:} plan risk mgmt (roles, scales, thresholds, reserves, cadence) \ra\ \lbl{identify} (brainstorming, \lbl{Delphi} = anonymous expert rounds, interviews, SWOT, checklists, RBS) \ra\ \lbl{qualitative} (probability $\times$ impact matrix \ra\ High/Med/Low priority) \ra\ \lbl{quantitative} (EMV, decision tree, Monte Carlo = thousands of simulated runs \ra\ probability of meeting date/budget) \ra\ plan \& implement responses \ra\ monitor (triggers, residual \& secondary risks, new risks). \textbullet\ \concept{Threat responses:} \kh{Avoid} (change plan/scope to eliminate), \kh{Mitigate} (reduce P or I --- extra testing, pilot, training), \kh{Transfer} (insurance, fixed-price contract, vendor SLA penalties --- risk still exists, owner changes), \kh{Accept} (active = contingency reserve/fallback; passive), \kh{Escalate} (outside project authority \ra\ program/portfolio/exec, e.g.\ company-wide cyber, regulatory change). \concept{Opportunities:} exploit (make it happen), enhance (raise P/I), share (partner), accept, escalate. \textbullet\ \concept{EMV} = \kh{probability $\times$ impact}; decision tree: choose option with best EMV (lowest expected cost). \lbl{Practice Q3b:} no extra test = $0.25\times120\text{k}=$ \kh{\$30k}; with test = $24\text{k}+0.10\times120\text{k}=$ \kh{\$36k} \ra\ testing costs \$6k more in expectation \ra\ \lbl{do not approve on EMV alone} (unless risk-averse: reputational/regulatory impact, \$120k worst case). \lbl{QB Q25:} vendor outage $0.3\times200\text{k}=60\text{k}$ (mitigate/transfer: SLA, rollback, cutover rehearsal), engineer $0.5\times40\text{k}=20\text{k}$ (mitigate: knowledge transfer, cross-train, retention), regulation $0.1\times500\text{k}=50\text{k}$ (monitor/escalate, design for compliance) \ra\ rank 60k $>$ 50k $>$ 20k; \lbl{contingency} $\approx$ sum EMV (130k). Limits: P \& I are guesses, average not worst case, ignores risk appetite \& interdependence, low-P high-I events underweighted. \textbullet\ \concept{Risk register:} ID, description (cause-event-effect), category, P, I, score/priority, \lbl{owner}, response strategy \& actions, trigger, contingency/fallback, residual risk, status. Living doc. \textbullet\ \concept{Communication:} tailor --- execs summaries, team detailed logs, users practical guidance; transparent (NHS NPfIT \textsterling12bn abandoned 2011: risks hidden, clinicians not engaged). \concept{Agile risk:} reviews each sprint/retro, risky items early (fail fast), risk burn-down, shared ownership; hybrid = register + adaptive. \textbullet\ \lbl{Practice Q2c:} incomplete payment integration testing \ra\ \lbl{mitigate}: complete end-to-end tests before go-live, rollback plan; or \lbl{avoid}: delay launch / launch without online payment. Half of staff untrained \ra\ \lbl{mitigate}: extra/short sessions, job aids, floor-walkers/super-users, prioritise front-line staff before peak.\par

\par\smallskip\noindent\hlhighlight{W8 · Procurement \& outsourcing}\ \concept{3 processes:} \lbl{plan} (what/how/when/who; make-or-buy; SOW; contract type; criteria) \ra\ \lbl{conduct} (RFP/RFQ out, bidder conference, proposals, evaluate, negotiate, award) \ra\ \lbl{control} (monitor performance/SLAs, formal change, \kh{verify/accept before payment}, disputes, close contract). \textbullet\ \concept{Make-or-buy factors:} cost (upfront vs recurring licence, TCO), time-to-market, quality, expertise/skills, control \& confidentiality/IP, customisation/flexibility, strategic importance/core competence, vendor lock-in \& support, risk. \lbl{Practice Q3c:} make = full customisation \& control but limited experience \ra\ schedule/quality risk; buy = sooner, proven, vendor support but licence fees + less flexibility + lock-in. Typical rec: \kh{buy} if non-core and deadline matters (negotiate config/API, exit clauses) --- state assumptions. \textbullet\ \concept{SOW} = description of work, deliverables, timeline, standards, \lbl{acceptance criteria} (legal reference, reduces scope creep). \concept{RFP} = complex need, vendors propose solutions (eval: quality, approach, capability + price). \concept{RFQ} = well-defined standard items (tablets, licences), price/delivery driven. \textbullet\ \concept{Contracts:} \kh{Fixed price} (FFP, FPIF, FP-EPA) --- buyer risk low, seller high; needs clear scope; changes expensive. \kh{Cost-reimbursable} (CPFF, CPIF, CPAF) --- buyer risk high; R\&D/uncertain. \kh{Time \& materials} --- medium; staff augmentation, uncertain integration (use caps, reviews). \textbullet\ \concept{Vendor evaluation:} \lbl{mandatory pass/fail checklist first} (certs, insurance, security --- catering case = checklist), then \lbl{weighted scoring} ($\Sigma$ score$\times$weight; e.g.\ A 3.95 vs B 3.90 --- check evidence/references when close), \lbl{scorecard} for ongoing KPIs (on-time \%, defect rate). \textbullet\ \concept{Supplier dispute (QB Q6.2):} refer to contract/SOW acceptance criteria (97\% pass) \ra\ \lbl{withhold payment} for non-conforming items / pay only accepted; issue formal notice + defect evidence; require rework at supplier cost (in-scope defect, not a change); escalate per contract; keep relationship via joint review meeting, recovery plan, scorecard. New feature request = \lbl{change control} (assess, approve, update contract/SOW) --- different from fixing a defect. \textbullet\ \concept{Outsourcing:} complete/partial, BPO, IT; onshore, nearshore, offshore, hybrid, multisource; insourcing = opposite. + cost, expertise, speed, focus on core; $-$ communication/time zones, quality control, confidentiality, dependency, hidden costs. Governance: strategic alignment, contract mgmt, roles, escalation, KPIs. Procurement risks: cost, schedule, quality, legal/IP, operational \ra\ fixed price, liquidated damages, strict acceptance criteria, backup supplier, knowledge transfer \& code/doc access.\par

\par\smallskip\noindent\hlhighlight{W9 · Leadership, communication, governance, alignment}\ \concept{Leadership} (people, vision, inspire, long-term) vs \concept{management} (tasks, plan, control, short-term efficiency). Styles: transformational, transactional, autocratic, laissez-faire, \lbl{adaptive/agile} (situational), servant (Scrum). Skills: communication, EI, negotiation, decision-making, stakeholder mgmt. \textbullet\ \concept{Communication mgmt (3):} plan (who needs what, channel, frequency, owner, 2-way feedback) \ra\ manage (distribute, store) \ra\ monitor (effectiveness, gaps, adjust). Channels $=n(n-1)/2$. Weekly report for routine; urgent security issue straight to sponsor + security lead. \textbullet\ \concept{Governance} = framework of policies, roles, decision rights, processes to direct \& control projects. \lbl{Structure:} \kh{Sponsor} (champions, funds, ultimate accountability, resolves escalations), \kh{Steering committee} (senior reps; approves charter/business case, major changes, monitors, resolves conflicts, confirms benefits), \kh{PMO} (standards, templates, KPIs, training, support). \lbl{Mechanisms:} stage/phase-gate reviews (go/kill), dashboards, escalation paths, decision frameworks, audits, change control board. Scale to project size. Challenges: unclear decision rights, stakeholders excluded, poor info, weak escalation. \textbullet\ \concept{Alignment} = project goals fit strategy: strategic, business-IT, stakeholder, team. \lbl{SAM} (Henderson \& Venkatraman): 4 domains business strategy, IT strategy, org infrastructure \& processes, IT infrastructure \& processes; strategic fit + functional integration. \lbl{Balanced Scorecard:} financial (ROI, cost), customer (satisfaction, adoption), internal process (cycle time, downtime, errors), learning \& growth (skills, training). \lbl{Value streams} (operational vs development). Also COBIT, ITIL, TOGAF. \textbullet\ \kh{Leadership drives people}, \kh{governance sets rules}, \kh{alignment picks the right game}; failures reinforce each other (poor leader \ra\ misalignment; weak governance \ra\ delayed decisions; misalignment \ra\ conflict).\par

\par\smallskip\noindent\hlhighlight{W10 · Integration \& change management}\ \concept{Integration} = ``glue'' coordinating all knowledge areas across lifecycle. \kh{7 processes:} develop charter; develop PM plan; direct \& manage work; manage knowledge (lessons learned); monitor \& control work; \lbl{perform integrated change control}; close project/phase. PMBOK = 5 process groups $\times$ 10 knowledge areas (Integration, Scope, Schedule, Cost, Quality, Resource, Comms, Risk, Procurement, Stakeholder). \textbullet\ \concept{Integrated change control:} CR logged \ra\ assess impact across scope/time/cost/quality/risk/resources \ra\ options (defer to release 2, trade-off, extra funding) \ra\ \lbl{CCB}/sponsor decision \ra\ update baselines \& plans \ra\ communicate. Security fix 2 wks before launch \ra\ \kh{assess impacts, compare options,} \kh{submit recommended change for approval} (not unilateral). \textbullet\ \concept{Closure:} formal acceptance, close contracts (verify obligations, final payment), release resources, lessons learned, archive, \lbl{handover to operations/ITSM}, final report, benefits tracking. \textbullet\ \concept{Change management} (people side): project = ``did we deliver it well?''; change = ``\kh{are people using it?}''. Success = delivery + adoption. Kübler-Ross curve (shock/denial \ra\ anger \ra\ bargaining \ra\ low \ra\ exploration \ra\ acceptance); ADKAR (Awareness, Desire, Knowledge, Ability, Reinforcement); Kotter. Actions: communicate the why, map affected stakeholders, champions, training \& job aids, pilots, support after go-live, quick wins, \lbl{measure uptake} (usage, feedback). IT project traits: complex, fast-changing tech, intangible, many stakeholders, uncertain, sensitive to requirements.\par

\par\smallskip\noindent\hlhighlight{W11--12 · ITSM \& ITIL 4}\ \concept{ITSM} = policies/processes to design, deliver, manage, improve IT services. \concept{Service} = \kh{means of enabling value co-creation} by facilitating outcomes customers want, \kh{without them managing specific costs and risks}. Traits: intangible, outcome-focused, co-created, continuous, risk/cost transfer. Product = owned/licensed, versioned. \lbl{PM vs ITSM:} temporary, unique, create/change, PM/sponsor vs ongoing, reliable delivery, service manager/service desk; link via project-to-service transition, shared change/config/knowledge mgmt, outputs vs outcomes, feedback loop. \textbullet\ \concept{Relationship models:} B2B (efficiency, reliability, SLAs --- supplier\ra council), B2C (UX, convenience), B2G (compliance, long contracts), G2B, \kh{G2C} (council\ra residents: accessibility, accuracy, trust), G2G, C2C, C2B, B2E, D2C, B2M, M2M (IoT), N2C, N2B. \textbullet\ \concept{4 dimensions} (holistic): \lbl{Organisations \& people} (culture, roles, skills --- untrained staff), \lbl{Information \& technology} (data, tools, knowledge base --- inaccurate data, disconnected systems), \lbl{Partners \& suppliers} (contracts, vendor data quality --- supplier SLAs), \lbl{Value streams \& processes} (workflows --- handoffs, procedures). \textbullet\ \concept{SVS:} opportunity/demand \ra\ [guiding principles, governance, service value chain, practices, continual improvement] \ra\ value. \kh{SVS = whole system;} \kh{SVC = operating model inside it.} \textbullet\ \concept{7 guiding principles:} focus on value; \kh{start where you are} (assess/reuse existing before replacing --- fix booking procedures before buying new system); \kh{progress iteratively with feedback} (small steps, pilot, measure); collaborate \& promote visibility; think \& work holistically; keep it simple \& practical; optimise \& automate (optimise before automating). \textbullet\ \concept{SVC 6 activities:} \lbl{Plan} (vision, policies, priorities, targets), \lbl{Improve}, \lbl{Engage} (stakeholders, demand, feedback), \lbl{Design \& Transition} (new/changed services, procedures, training, into live), \lbl{Obtain/Build} (components, tools, suppliers), \lbl{Deliver \& Support} (operate, service desk). Value streams combine activities per outcome. \textbullet\ \concept{34 practices:} 14 general (continual improvement, info security, risk, knowledge, project, relationship...), 17 service (incident, problem, change enablement, service request, service desk, SLM, IT asset, config, release, availability, capacity, continuity), 3 technical (deployment, infra \& platform, software dev). \textbullet\ \concept{Incident mgmt} = \kh{restore normal service ASAP} (identify, log, categorise, prioritise by impact $\times$ urgency, diagnose, resolve/workaround, close). \concept{Problem mgmt} = \kh{find \& remove root cause} of (recurring) incidents: detect, RCA, \lbl{workaround}, \lbl{known error} record, permanent fix via change. \concept{Change enablement:} RFC, assess risk/impact, approve (CAB), schedule, implement, review; standard/normal/emergency changes. \concept{Service desk} = SPOC, logs, fulfils requests, escalates, communicates, updates knowledge base. \concept{Knowledge mgmt:} consistent guidance, self-service FAQs, single source of truth. \textbullet\ \concept{SLM:} define \& agree \lbl{SLAs} with measurable targets, monitor \& report, regular service reviews, improve. Good targets: \kh{availability \%} (e.g.\ 99.9\% during store hours), \kh{response time} (95\% of stock queries $\le$2 s), \kh{incident response/resolution by priority} (P1 $\le$15 min / 4 h), \kh{data freshness} (stock sync $\le$5 min), \kh{CSAT $\ge$4/5}, first-contact resolution \%. Monitor via automated dashboards, monthly reports, quarterly reviews, breach escalation \& service credits. ``Reliable and responsive'' is not measurable. Avoid \lbl{watermelon SLA} (green metrics, unhappy users) --- add experience measures. \textbullet\ \concept{Availability} $=\dfrac{\text{agreed time}-\text{downtime}}{\text{agreed time}}\times100\%$. \lbl{QB Q28:} $(720-6)/720=$ \kh{99.17\% $<$ 99.5\%} \ra\ SLA breached (allowed $0.5\%\times720=3.6$ h). Availability alone misses queue/wait time, usability, support quality \ra\ add job turnaround (95\% start $\le$30 min), consistent-guidance/first-contact resolution, CSAT. \textbullet\ \concept{Continual improvement model:} What is the vision? \ra\ Where are we now? (baseline metrics) \ra\ Where do we want to be? (targets) \ra\ How do we get there? \ra\ Take action \ra\ Did we get there? \ra\ How do we keep momentum? Use data + user feedback before big infrastructure spend (fix scheduling \& instructions first, then measure). \textbullet\ \concept{ITIL adoption challenges:} cultural resistance, trying all 34 practices at once, lack of executive support. Automation/AI in ITSM: ticket routing, chatbots/self-service, predictive monitoring, AI-assisted change risk assessment. \textbullet\ \lbl{ITSM quick distinctions:} \kh{incident} = unplanned interruption/degradation (restore fast) vs \kh{service request} = normal ask (password reset, new laptop; fulfil via catalogue) vs \kh{problem} = cause of one or more incidents vs \kh{known error} = problem with documented root cause/workaround. \lbl{Priority} = impact (users/business affected) $\times$ urgency (how fast it hurts): P1 whole LMS down in exam week, P2 one faculty, P3 single user with workaround, P4 cosmetic. \lbl{Change types:} \kh{standard} (pre-authorised, low risk, routine: add user, patch template), \kh{normal} (assessed \& authorised by change authority/CAB), \kh{emergency} (urgent fix, expedited approval, review afterwards). \lbl{SLA} (provider \ensuremath{\leftrightarrow} customer) vs \lbl{OLA} (internal teams supporting SLA) vs \lbl{underpinning contract} (with external supplier); SLA targets must be backed by OLAs/UCs. \lbl{Project \ra\ service transition:} handover docs, KB articles, trained service desk, support model \& SLAs agreed, early-life support, warranty period, acceptance into operations --- otherwise ``thrown over the wall''. \textbullet\ \lbl{EVM interpretation phrases:} SPI 0.80 = ``achieving 80\% of planned progress'', CPI 0.75 = ``earning \$0.75 of work per \$1 spent'', EAC = ``if trend continues, final cost \$X, overrun \$VAC''; always add cause + action + whether compression worsens CPI. \lbl{Float phrase:} ``C has 5 d float, so a 3 d delay is absorbed; end date unchanged; C's path now has 2 d float left''.\par

\par\smallskip\noindent\hlhighlight{Question bank · W7--W12 answer skeletons}\ \lbl{QB Q6.1 museum batch:} pass $=(1750-105)/1750=1645/1750=$ \kh{94.0\% $<$ 97\%} \ra\ \kh{batch rejected} (allowed failures $3\%\times1750=52.5$, i.e.\ $\le$52). \lbl{QA} (prevention, process): agree scanning standard + calibration checklist + pilot batch/sample sign-off with supplier before full run \ra\ stops defects being produced. \lbl{QC} (detection, product): inspect every batch against resolution/readability criteria, log failures by cause (Pareto), return non-conforming images \ra\ only conforming output accepted. \textbullet\ \lbl{Q6.3 publish everything now?} Museum = \kh{low risk appetite} (public trust, legal/permission \& privacy exposure, reputational damage, takedown cost) vs small visitor gain \ra\ sponsor's plan exceeds appetite. Response: \kh{avoid} for affected items (withhold until rights/sensitivity cleared) + \lbl{mitigate} (permissions review workflow, sensitivity checklist, takedown process) while publishing cleared photos now (phased launch keeps visitor benefit). Governance: escalate to steering committee/sponsor with risk register entry, options \& EMV-style impacts; documented decision with owner. Stakeholders: consult curators, legal/rights holders, community groups; communicate phased plan publicly. \textbullet\ \lbl{QB Q26 smart waste:} (a) supplier \ra\ council = \kh{B2G} priority = data accuracy/reliability per SLA; (b) council \ra\ residents = \kh{G2C} priority = accurate, timely, accessible notifications \& trust. 4D: people (untrained service staff \ra\ training + scripts + knowledge base), info \& tech (inaccurate notifications, disconnected systems \ra\ integrate tracking with notification system, data validation), partners (supplier data quality \ra\ SLA with data-accuracy targets \& review), value streams (missed-pickup report \ra\ crew dispatch handoffs \ra\ end-to-end workflow, single ticket). Availability $\ne$ value: system up but wrong notifications = no outcome; measure \% notifications accurate ($\ge$98\%), missed pickups resolved $\le$48 h, CSAT. \textbullet\ \lbl{QB Q27 double bookings:} start where you are = assess current procedures/data, reuse what works, find whether cause is procedure or software before replacing; iterative = fix booking rules \& staff procedures, pilot one site, measure double-booking rate, then decide on tech. SVC: \lbl{Plan} (target zero double bookings, policy), \lbl{Design \& Transition} (redesigned booking procedure, config fix, training), \lbl{Improve} (track incidents, feedback, iterate) --- linked by measurement feeding next plan. SVS: tech alone fails without governance (decision rights, policy), practices (incident/problem/change, knowledge, SLM) and value chain working together around customer value. \textbullet\ \lbl{QB Q28.3 before buying compute:} (1) \kh{knowledge mgmt}: single job-submission guide, FAQs, templates, consistent support scripts, KB articles from common tickets; (2) \kh{continual improvement}: baseline queue time \& failed submissions, optimise scheduler (priority queues, fair-share, off-peak batching), then re-measure. Evaluate with data (queue wait, failed/resubmitted jobs, first-contact resolution, ticket volume) + researcher feedback (survey/CSAT) before vs after; buy capacity only if targets still missed.\par

\par\smallskip\noindent\hlhighlight{Practice exam (4 Q $\times$ 15 = 60 pts) · full answers}\ \lbl{1a Uni app: Waterfall vs Agile (5).} \concept{Planning:} Waterfall = detailed upfront plan, fixed scope/phases; Agile = high-level roadmap, detailed planning each sprint (rolling). \concept{Stakeholders:} Waterfall = mainly at requirements sign-off and final acceptance; Agile = continuous, sprint reviews/demos, PO represents students. \concept{Change:} Waterfall = formal CR, costly late; Agile = welcomed, reprioritise backlog each sprint. \kh{Recommend Agile/Scrum}: must release core (enrolment, announcements, digital ID) before semester = MVP in first releases; features keep coming through the year = incremental releases; student feedback shapes later features. Keep privacy/security requirements and release milestones fixed (light governance). \textbullet\ \lbl{1b Scope creep (5).} It is scope creep because event registration + navigation are \kh{added without impact assessment or approval} while \kh{deadline and budget unchanged} (baseline not updated). Consequences (any 2): missed release before semester; budget overrun; rushed/less testing \ra\ defects in core enrolment/ID; team overtime/burnout; stakeholder conflict over priorities. Measure: \kh{integrated change control} --- every request logged as CR, PM assesses impact on scope/time/cost/quality/risk, sponsor/CCB approves, defers or rejects; approved items go to product backlog for a later release, baseline/plan updated; useful changes still considered without derailing release 1. \textbullet\ \lbl{1c PERT O=4 M=7 P=16 (5).} $t_E=(4+4\times7+16)/6=48/6=$ \kh{8 days}; $\sigma=(16-4)/6=2$ d. Useful: combines best/likely/worst case, weights most likely, gives realistic estimate + measure of uncertainty for buffers and risk. Limitation: 8 d is an \emph{average} --- skewed by pessimistic 16 d, only $\approx$50\% chance of finishing $\le$8 d, could take up to 16 d; estimates are subjective guesses \ra\ committing to 8 d is risky; commit $t_E+\sigma$ to $+2\sigma$ (10--12 d) or add buffer. \textbullet\ \lbl{2a Requirements (4).} Important: aligns system with resident \& staff needs; reconciles different stakeholder expectations; reduces scope creep, rework and defects; basis for scope, design, estimates, testing/acceptance. Techniques (2, justified): \kh{workshops} with planning staff, assessors, residents' reps --- surface \& resolve conflicting expectations together, agree priorities; \kh{interviews} with permit officers (detail of rules/documents); or \kh{observation} of current permit process; \kh{surveys} of many residents; \kh{prototypes} of the online form to confirm usability. \textbullet\ \lbl{2b Leadership (5).} One directive style suits new council staff (need clear instructions, structure) but \kh{demotivates experienced developers} (no autonomy, slows decisions, stifles creativity; Herzberg: responsibility/autonomy are motivators) \ra\ unsuitable as one-size. Adapt (\kh{situational/adaptive}): experienced devs \ra\ delegating/participative, set goals \& standards, let them choose how, consult on technical decisions \ra\ ownership, speed, engagement; new staff \ra\ directive + coaching: clear tasks, checklists, training, mentoring by seniors, frequent check-ins, gradually more autonomy as competence grows \ra\ fewer errors, confidence, skill growth. Result: higher morale, productivity, better collaboration. \textbullet\ \lbl{2c Risks before busy period (6).} \lbl{R1 incomplete payment integration testing} \ra\ payments fail/duplicate at peak, applications lost, revenue \& reputation damage, emergency fixes. Response \kh{mitigate}: complete end-to-end integration testing with payment provider before go-live (test cases, load at peak volume), rollback/manual-payment fallback; or \kh{avoid}: delay launch / launch without online payment. \lbl{R2 only half staff trained} \ra\ staff cannot process/assist applications, errors, backlogs, poor resident experience in busiest period. Response \kh{mitigate}: schedule extra short sessions for remaining staff before launch, quick-reference guides, super-users/floor-walkers on go-live; owner = training lead. \textbullet\ \lbl{3a WBS (4).} Purpose: \kh{deliverable-oriented hierarchical decomposition of 100\% of scope} into work packages that can be estimated, scheduled, assigned. Scenario WBS: Training platform \ra\ Course module, Assessment engine, Progress/PD tracking, User mgmt \& integration (HR/SSO), Testing, Training \& deployment, PM. Supports planning: work packages \ra\ bottom-up cost estimates \& cost baseline; activities \& dependencies \ra\ schedule; owners (RACI) \ra\ accountability; nothing missed or double counted \ra\ scope control; basis for risk identification. \textbullet\ \lbl{3b EMV (5).} No testing: $0.25\times120{,}000=$ \kh{\$30,000}. With testing: $24{,}000+0.10\times120{,}000=24{,}000+12{,}000=$ \kh{\$36,000}. Testing raises expected cost by \$6,000 \ra\ \kh{do not approve on EMV alone}. Caveat: if org is risk-averse (assessment defects hit staff certification/reputation), could still justify; or seek cheaper targeted testing. \textbullet\ \lbl{3c Make-or-buy (6).} \concept{Cost:} make = dev effort, no licence, maintenance in-house; buy = lower upfront, recurring licence fees (TCO over years). \concept{Time:} make = slower, learning curve (limited experience) \ra\ schedule risk; buy = available sooner. \concept{Expertise/quality:} make = limited experience \ra\ defect risk; buy = proven, vendor support. (Alt: \concept{flexibility/control}: make = full customisation \& IP; buy = less flexible, vendor lock-in.) \kh{Recommend buy} (assessment authoring is not core differentiator, speed matters, team inexperienced) --- negotiate configuration/API access, SLA, exit clause. Assumptions: licence cost within budget over 3--5 yrs, component meets core requirements, deadline important. \textbullet\ \lbl{4a EVM (4).} $SPI=EV/PV=240k/300k=$ \kh{0.80} \ra\ behind schedule, only 80\% of planned work done. $CPI=EV/AC=240k/320k=$ \kh{0.75} \ra\ over budget, \$0.75 of value per \$1 spent. Both $<1$ \ra\ behind and over budget; (SV $-$60k, CV $-$80k); if trend continues EAC = BAC/0.75, investigate causes, corrective action. \textbullet\ \lbl{4b Tuckman (5).} Forming \ra\ storming \ra\ norming \ra\ performing \ra\ adjourning. Team was \kh{performing}; new members change dynamics \ra\ team \kh{regresses to forming/storming} (roles, norms re-negotiated, disagreements on practices, slower delivery). Sponsor wrong: experience does not prevent regression; development is cyclical when membership changes. Intervention: \kh{team working-agreement/charter workshop} to revisit norms, practices, roles (+ buddy/mentor onboarding, retrospective) \ra\ new members understand \& contribute to norms, conflict resolved openly \ra\ moves quickly back to norming/performing. \textbullet\ \lbl{4c SLM (6).} ``Reliable and responsive'' is not measurable. Targets: (1) \kh{availability $\ge$99.9\%} during store trading hours (6am--11pm) per month --- stores cannot track stock/sell without it; (2) \kh{stock level updates visible $\le$5 min} / 95\% of inventory queries respond $\le$2 s --- accurate real-time stock for replenishment \& purchasing; (alt) P1 incident response $\le$15 min, resolution $\le$4 h. Monitor \& review: agreed in SLA, automated monitoring dashboards, monthly SLA reports, quarterly service review meetings, breach escalation, service credits/penalties, improvement plans, review targets as business changes. Assumptions: trading hours, store count, peak periods, data volumes. \textbullet\ \lbl{Likely variants of the same questions:} \concept{methodology} --- regulated/fixed-price \ra\ Waterfall, mixed \ra\ hybrid (Waterfall governance gates + Agile build sprints); always compare planning, stakeholders, change, documentation. \concept{scope creep} --- also gold-plating (team adds unrequested features) \ra\ same control: CR + impact + approval + backlog. \concept{PERT} --- if asked probability: $t_E\pm1\sigma\approx68\%$, $\pm2\sigma\approx95\%$; path estimate = sum of $t_E$, $\sigma_{path}=\sqrt{\Sigma\sigma_i^2}$. \concept{requirements} --- add NFR rewrite (measurable, testable) + RTM. \concept{leadership} --- transformational for change/vision, servant for Scrum teams, transactional for routine targets. \concept{risk} --- opportunity version: exploit/enhance/share/accept; always risk + impact + strategy + concrete action + owner. \concept{EMV} --- if EMV(with action) $<$ EMV(without) \ra\ approve; mention risk appetite either way. \concept{make-or-buy} --- other factors: control/IP, flexibility, vendor lock-in, strategic/core, support, security of data. \concept{EVM} --- also CV, SV, EAC=BAC/CPI, VAC, ETC; CPI $<1$ \& SPI $<1$ \ra\ root cause, reforecast, scope triage; avoid crashing when CPI already low. \concept{Tuckman} --- other triggers for regression: new PM, new goals, re-org, remote members; interventions: team charter, RACI, retrospectives, team-building, coaching. \concept{SLM} --- other targets: P1 response $\le$15 min/resolve $\le$4 h, first-contact resolution $\ge$70\%, CSAT $\ge$4/5, report generation $\le$10 s, data accuracy $\ge$99.5\%; make each SMART + why it matters + how measured.\par

\end{multicols}
\end{document}
